\section{Walk the complete eighty-two-byte example}
The byte offsets below are relative to the beginning of the file. Sizes in this table
are numerical decimal values; the identifiers are literal ASCII bytes. The data region
begins at 56, after two complete earlier chunks and the data chunk's own header.

\tagpdfsetup{table/header-rows={1}}
\begin{center}
  \begin{tabular}{rll}
    Offset & Region                                    & Extent in bytes \\
    \hline
    0      & RIFF identifier, size 74, WAVE identifier & 12              \\
    12     & JUNK header, length three                 & 8               \\
    20     & JUNK payload and its pad                  & $3+1$           \\
    24     & fmt header, length sixteen                & 8               \\
    32     & PCM format payload                        & 16              \\
    48     & data header, length sixteen               & 8               \\
    56     & Four stereo frames                        & 16              \\
    72     & note header, length one                   & 8               \\
    80     & note payload and its pad                  & $1+1$           \\
  \end{tabular}
\end{center}

The final pad is \texttt{7f}. The chunk walk ends at 82, exactly at the RIFF limit.
Finding the data chunk at 48 is not permission to stop validating: its following chunk
must still fit. A file with the same sixteen sample bytes and a damaged final chunk
must fail parsing without publishing the descriptor.

Record the current offset and remaining extent before each header load. Then annotate
the one place where the next offset gains an extra byte for odd size. This trace makes
an omitted pad visible as a wrong next identifier rather than as an unexplained sample
error much later in the program.
